% AMS-LaTeX conversion of the 1998 manuscript (original file: howe6.tex).
% Original mathematical text, theorem labels, and bibliography are retained.
\documentclass[12pt,letterpaper]{amsart}
\usepackage[T1]{fontenc}
\usepackage{amsmath,amssymb}
\usepackage{newtxtext,newtxmath} % Times-style text and matching math fonts
\usepackage{microtype}
     % Please give any \input statements here:

%     Author: Please fill in the information for the following macros.
%     Directions given after the comment characters should be deleted.

\title[The generalized Hecke algebra of
$\overline{SL(2,F)}$]{A correspondence for the generalized Hecke algebra of
the metaplectic cover $\overline{SL(2,F)}$, $F$ p-adic} 
% If the title is too long for the running heads, give a
         % short version as an optional argument, with the syntax
         % \title[Short Title]{Full Title}. 
                          % (The syntax \title{}[], which works in
                          % some styles, does not work in this one.)

  % Give separate author, address, and email entries for each author.
  % The \author entry for the second author should follow the \email
  %     for the first author, etc. 
\author{David Joyner} 
 % A shortened version for the running heads may be given in []'s,
           % using the same syntax as for titles.
\address{Mathematics Department, Chauvenet Hall,
572 Holloway Road, Annapolis, MD 21402, USA} 

\email{wdj@nadn.navy.mil}
 % You will need @@ for an "at" sign in amslatex.

           % If the combined shortened forms of the authors' names is
           % too long for the running heads, you must set the
           % running heads explicitly by giving these commands:


% \pagestyle{myheadings} \markboth{<name of first author>, et al.}{<title>}



          % \thanks entries are to acknowledge grants. You may combine
          % all acknowledgments into one \thanks entry, or may use
          % multiple \thanks entries. They generate footnotes without
          % tags, so you must be explicit about which authors are
          % thanking whom.
\thanks{I thank Bill Wardlaw for some helpful 
conversations on the proof of Proposition \ref{prop:main}
and the anonomous referee for many helpful corrections
and suggestions.}
   % optional

       % Keywords are important, as they may be used by browsers to
       % search the archive for papers of interest. So think about
       % how people might look for your paper.
\keywords{metaplectic groups, generalized Hecke algebras,
representation theory, p-adic groups, SL(2)}

       % Give the appropriate broad classifications from the AMS
       % subject classifications. For instance, group actions papers 
       % should include a reference here to 57S:
\subjclass[1991]{22E,11F}

       % Abstracts are required, and will be distributed by a listserv list.
       % Please place the text of your abstract in the argument of the
       % following definition.

\newcommand{\abstracttext}{We prove, using a technique developed for 
$GL(n)$ in Howe and Moy \cite{H}, 
a bijection between generalized Hecke
algebras of $G=SL(2,F)$ over a $p$-adic field
and those of its $n$-fold metaplectic cover $\overline{G}$. 
This result implies that
there is a canonical correspondence
between irreducible admissible representations of $G$ and 
genuine irreducible admissible representations of $\overline{G}$
of ``sufficiently large level" (depending on $n,p$).}

% Here are some recommended declarations for theorem styles. They use
% the convention that all theorem-like environments are numbered
% consecutively by section.

\newtheorem{thm}{Theorem}[section]
\newtheorem{prop}[thm]{Proposition}
\newtheorem{lem}[thm]{Lemma}
\newtheorem{cor}[thm]{Corollary}

\theoremstyle{definition}

\newtheorem{defn}[thm]{Definition}
\newtheorem{notation}[thm]{Notation}
%\newtheorem{example}[thm]{Example}
\newtheorem{conj}[thm]{Conjecture}
\newtheorem{prob}[thm]{Problem}

\theoremstyle{remark}

\newtheorem{rem}[thm]{Remark}


% Place any personal macros used in the document here.

\newtheorem{theorem}{Theorem}
\newtheorem{corollary}[theorem]{Corollary}
\newtheorem{conjecture}[theorem]{Conjecture}
\newtheorem{lemma}[theorem]{Lemma}
\newtheorem{proposition}[theorem]{Proposition}
\newtheorem{definition}[theorem]{Definition}
\newtheorem{example}[theorem]{Example}
\newtheorem{axiom}{Axiom}
\newtheorem{remark}{Remark}
\newtheorem{exercise}{Exercise}[section]


\begin{document}
\begin{abstract}
\abstracttext
\end{abstract}
\maketitle

%%%% **** The text of the paper starts here **** %%%%

Let $F$ be a p-adic field of characteristic 0, 
with uniformizer $\pi$, ring of integers $\mathcal{O}$, and let 
$q$ denote the cardinality of the residue field.
Let $\mu_n(F)$ denote the group of $n^{th}$ roots of unity
in $F$ and assume $|\mu_n(F)|=n$. We 
will identify $\mu_n(F)$ with a subgroup 
of ${\mathbb{C}}^\times$
(via some fixed isomorphism 
$\theta:\mu_n(F)\rightarrow \mu_n({\mathbb{C}})$).
Write $\mu_n=\mu_n(F)$.
Let $G=SL(2,F)$ and let $\mathfrak{g}$ denotes its Lie algebra.
Let $\overline{G}=\overline{SL(2,F)}$ denote the n-fold metaplectic
cover defined by the cocycle
\[
\beta(g_1,g_2)=({x(g_1g_2)\over x(g_1)},{x(g_1g_2)\over x(g_2)}),
\]
for $g_1,g_2\in G$, where
\[
x(
\begin{array}{cc}
a & b\\
c & d
\end{array}
)=
\left\{
\begin{array}{c}
c,\ \ \ \ \ {\rm if}\ c\not= 0,\\
d,\ \ \ \ \ {\rm otherwise},
\end{array}
\right.
\]
and where $(\ ,\ )$ is the $n^{th}$ power Hilbert symbol. 
By \S 1 of Kubota \cite{K}, $\beta$ is a 2-cocycle 
(called a ``factor set'' in \cite{K}) which is trivial
near the identity in $G\times G$, so $\overline{G}$ is a 
topological covering group of $G$. We denote
elements of $\overline{G}$ by $\overline{x}=(x,\zeta)$,
$x\in G,\ \zeta\in \mu_n$, and multiplication by
\[
\overline{x}_1\overline{x}_2=(x_1,\zeta_1)(x_2,\zeta_2)
=(x_1x_2,\beta(x_1,x_2)\zeta_1\zeta_2).
\]
Let $\rho:\overline{G}\rightarrow G$ denote
the natural surjection. If $H$ is any subgroup
of $G$ then we define $\overline{H}=\rho^{-1}(H)$.
We choose a Haar measure on $G$, to be normalized
below, and a Haar measure on $\mu_n$ normalized so that 
$vol(\mu_n)=1$. Let the Haar measure on $\overline{G}$ be
the product measure.

We prove, using a technique developed for $GL(n)$ in \cite{H}, 
a bijection between generalized Hecke
algebras of $G$ and those of $\overline{G}$. 
A precise statement will be
given in Theorem \ref{thrm:main}. This result implies that
there is a canonical one-to-one correspondence
between irreducible admissible representations of $G$ and 
genuine irreducible admissible representations of $\overline{G}$
of ``sufficiently large level" (depending on $n,p$).
For the precise statement, see Corollary
\ref{corollary:main}. 

The study of genuine admissible representations
on p-adic groups has its origins in the
book \cite{G} by S. Gelbart. Flicker \cite{F}
determined a comparison of Hecke algebras of
{\it smooth functions} on $GL(2,F)$ and 
$\overline{GL(2,F)}$ using the matching of orbital integrals.
This was then used, with the Arthur-Selberg trace formula, to 
establish a global ``Shimura correspondence''.
Neither of these two works used the Hecke algebras considered here.
G. Savin \cite{S} has given an isomorphism between
the Iwahori-Hecke algebras in the case when $p$ is
relatively prime to $2n$ and $G$ is any reductive
$p$-adic group which splits over $\mathcal{O}$. This isomorphism 
generalized a result of Flicker and Kazhdan \cite{FK} who
considered the special case of $GL(r)$. 
(Savin's result is discussed briefly
in \S \ref{section:savin} below.)
The works Rallis-Schiffmann \cite{RS} and Schultz
\cite{Sch} also dealt with a correspondence between
representations on $\overline{G}$ and on $G$ (or a closely related 
group),
though from a different perspective and using different methods.
Our hope is that the result in the present paper will form
one small step, along with Waldspurger's
multiplicity one theorem for $\overline{SL(2,F)}$ 
\cite{W} and the Arthur-Selberg trace formula, in an eventual proof 
Labesse's muliplicity one conjecture for $SL(2)$. 

\section{Some notation}

Let
\[
s_1=
\left(
\begin{array}{cc}
0 & -1\\
1 & 0
\end{array}
\right),\ \ \ \ 
s_2=
\left(
\begin{array}{cc}
0 & -\pi^{-1}\\
\pi & 0
\end{array}
\right),\ \ \ \ 
t=
\left(
\begin{array}{cc}
0 & -1\\
\pi & 0
\end{array}
\right),\ \ \ \ 
t_0=
\left(
\begin{array}{cc}
\pi & 0 \\
0 & \pi^{-1}
\end{array}
\right).
\]
Note $t_0=s_1^{-1}s_2$, $t\notin G$ and $ts_1t^{-1}=s_2$. 

Let
\[
u_+(x)=\left(
\begin{array}{cc}
1 & x\\
0 & 1
\end{array}
\right),\ \ \ \ 
u_-(x)=\left(
\begin{array}{cc}
1 & 0\\
x & 1
\end{array}
\right),\ \ \ \ 
h(x)=\left(
\begin{array}{cc}
x & 0\\
0 & x^{-1}
\end{array}
\right).
\]

Let
\[
B=\{
\left(
\begin{array}{cc}
\mathcal{O} & \mathcal{O}\\
\pi \mathcal{O} & \mathcal{O}
\end{array}
\right)\}\cap G
\]
denote the Iwawori subgroup,
\[
{\mathfrak{b}}=\{
\left(
\begin{array}{cc}
\mathcal{O} & \mathcal{O}\\
\pi \mathcal{O} & \mathcal{O}
\end{array}
\right)\}\cap \mathfrak{g}
\]
denote its Lie algebra, and let
\[
B_k=(I+t^{2k}\left(
\begin{array}{cc}
\mathcal{O} & \mathcal{O}\\
\pi \mathcal{O} & \mathcal{O}
\end{array}
\right))\cap G
\]
denote a $k^{th}$ filtration subgroup. Let
\[
{\mathfrak{b}}_k 
=\{
\left(
\begin{array}{cc}
a & b\\
\pi c & -a
\end{array}
\right)
\ |\ 
a,b,c\in \pi^{k}\mathcal{O}\}
\]
denote the Lie algebra of $B_k$. 

\begin{definition}
Let $H$ denote a locally compact unimodular group and 
let $K$ denote
an open compact subgroup of $H$ with Haar measure on 
$H$ normalized so that
$K$ has measure one. Let
$\mathcal{H}(H//K)$ denote the space of functions 
on $H$ which
are locally constant, compactly supported, and 
bi-$K$-invariant. This forms an
algebra under convolution,
\[
(f*g)(x)=\int_H f(y)g(y^{-1}x)dy,
\]
called the {\bf generalized Hecke algebra}.
\label{def:1}
\end{definition}

We shall study the algebra $\mathcal{H}(G//B_k)$ in section \ref{section:sl2}.

Recall that a 2-cocycle $\alpha :H\times H \rightarrow \mu_n$
of a group $H$ is {\bf trivial} if there is a function
$s:H\rightarrow \mu_n$ (called a {\bf splitting})
such that $\alpha(h_1,h_2)=s(h_1)s(h_2)s(h_1h_2)^{-1}$,
for $h_1,h_2\in H$. In this case the cover of $H$
defined by $\alpha$ is said to {\bf split}.

Recall that a complex-valued function $f$ on
$\overline{G}$
is called {\bf genuine} if $f(g,\zeta)
=\theta(\zeta) f(g,1)$, for
$g\in G$ and $\zeta\in \mu_n$. Here
$\theta$ is a fixed isomorphism 
$\theta:\mu_n\rightarrow \mu_n({\mathbb{C}})$.
It being agreed that we may identify
$\mu_n$ and $\mu_n({\mathbb{C}})$, we shall, {\it from this point on, 
drop the $\theta$ from the notation.}

\begin{definition} Assume that $k>0$ is chosen so large that
$\overline{B}_k$ splits. Assume that the Haar measure on
$\overline{G}$ is normalized so that $\overline{B_k}$ has 
measure 1. Let $\mathcal{H}(\overline{G}//B_k)_{gen}$
denote the space of functions on $\overline{G}$ which
are genuine, locally constant, compactly supported, and 
bi-$B_k$-invariant.
This is the {\bf (metaplectic) generalized Hecke algebra} 
on $\overline{G}$.
\end{definition}

We will study this algebra 
in section \ref{section:metaplectic}. Note that it
depends on the splitting of the cocycle on
$B_k$. To have a measure on $\overline{G}$ which fits
with the measure on $G$ used in \S \ref{section:sl2},
we shall also assume that $\mu_n$ has measure 1.

\section{Some lemmas}

Let ${\mathbb{F}}_q$ denote the residue field of $F$.

\begin{lemma} For $k>0$,
$B_k/B_{k+1}\cong {\mathfrak{b}}_k/{\mathfrak{b}}_{k+1}
\cong {\mathbb F}_q^3$, 
as abelian groups. (The second isomorphism is
not canonical.)
\end{lemma}

{\bf proof}:
The graph of the the first isomorphism is 
$\{(g\cdot B_{k+1},X+ {\mathfrak{b}}_{k+1})\ |\ 
(1+X+{\mathfrak{b}}_{k+1})\cap g\cdot B_{k+1}\not= 
\emptyset\}$. It is also a special case of
a general result of Morris (Proposition 3.4 in \cite{M}). 
The second isomorphism is due to the fact that
the map
\[
\left(
\begin{array}{cc}
a & b\\
\pi c & -a
\end{array}
\right)\longmapsto
\left(
\begin{array}{cc}
a\ ({\rm mod}\ \pi) & b\ ({\rm mod}\ \pi)\\
\pi (c\ ({\rm mod}\ \pi)) & -a \ ({\rm mod}\ \pi)
\end{array}
\right)
\]
has kernel $\cong {\mathbb F}_q^3$. $\Box$

The following result is stated as Proposition 2.2 in 
chapter 3 of \cite{H}. 

\begin{lemma} In the notation of Definition \ref{def:1},
let $f_h$ denote the characteristic function of the 
double coset $K h K$ in $H$. We have,
$f_1$ is a unit of $\mathcal{H}(H//K)$. More generally,
if
\[
vol(Kh_1K)vol(Kh_2K)=vol(Kh_1h_2K)
\]
then $f_{h_1}*f_{h_2}=f_{h_1h_2}$.
\label{lemma:general}
\end{lemma}

We need a version of the previous lemma for genuine 
functions. 

For $\overline{x}=(x,\zeta)\in\overline{G}$
and $g\in G$, let
\[
\overline{f}_{g}(\overline{x})
=
\left\{
\begin{array}{cc}
\zeta,& x\in B_kgB_k,\\
0,& {\rm otherwise}.
\end{array}
\right.
\]
Note that $\overline{f}_g$ is genuine. Let
\[
\phi_1*\phi_2(g)=
\int_{\overline{G}}\phi_1(x)\phi_2(x^{-1}g) dx,
\ \ \ g\in \overline{G}.
\]

Now we prove a genuine 
analog of the general Lemma \ref{lemma:general} above.

\begin{lemma} 
If $vol(B_k)=1$ and if 
\[
vol(B_kg_1B_k)
vol(B_kg_2B_k)
=vol(B_kg_1g_2B_k),\ \ \ \ {\rm for}\ g_1,g_2\in G,
\]
then $\overline{f}_{g_1}*\overline{f}_{g_2}=
\overline{f}_{g_1g_2}$.
\label{lemma:genuine}
\end{lemma}

proof: We compute
\[
\begin{array}{c}
\overline{f}_{g_1}*\overline{f}_{g_2}(g,1)=
\int_{\overline{G}}
\overline{f}_{g_1}(x)\overline{f}_{g_2}(x^{-1}(g,1))dx\\
=\int_{G}\overline{f}_{g_1}(x,1)\overline{f}_{g_2}((x^{-1}g,1))
\beta(x,x^{-1}g)^{-1}dx\\
=\int_{G}
char(\overline{B_k}(g_1,1)\overline{B_k})(x,1)
char(\overline{B_k}(g_2,1)\overline{B_k})(x^{-1}g,1)
\beta(x,x^{-1}g)^{-1}dx\\
=\int_{G}f_{g_1}(x,1)f_{g_2}(x^{-1}g,1)
\beta(x,x^{-1}g)^{-1}dx.
\end{array}
\]
The above equations imply that
the support of this integral is contained in the support
of $f_{g_1}*f_{g_2}$. By Lemma \ref{lemma:general},
$f_{g_1}*f_{g_2}=f_{g_1g_2}$, so 
\[
supp(\overline{f}_{g_1}*\overline{f}_{g_2})
\subset supp(\overline{f}_{g_1g_2}).
\]
Since the support each element of the Hecke algebra 
contains at least one double $B_k$-coset, the above
inclusion must be an equality. Therefore, we must have
$\overline{f}_{g_1}*\overline{f}_{g_2}=c\overline{f}_{g_1g_2}$
for some constant $c$. Integrating both sides over 
$G$ and using the hypothesis implies $c=1$.   $\Box$

As in the above proof, we have the following expression
for the convolution in terms of the cocycle.

\begin{lemma}
\label{lemma:convolve}
We have
\[
\begin{array}{c}
\overline{f}_{g_1}*\overline{f}_{g_2}(g,\zeta)
=\zeta\int_{G}\overline{f}_{g_1}(x,1)\overline{f}_{g_2}(x^{-1}g,1)
\beta(x,x^{-1}g)^{-1}dx,
\end{array}
\]
for any $g_1,g_2\in G$.
\end{lemma}

Note that the generalized metaplectic Hecke
algebra depends on the cocycle $\beta$
so let us denote this dependence temporarily by
$\mathcal{H}_\beta(\overline{G}//B_k)_{gen}$.
If the cocycle is changed to an equivalent one,
say to $\beta'$, where 
\[
\beta(g,h)=\beta'(g,h)s(g)s(h)s(gh)^{-1}
\]
is the cocycle modified by $s$ 
and $s:G\rightarrow \mu_n$ is any function
satisfying $s(1)=1$
(for example, the Kubota splitting
\cite{KP}), then the two algebras are 
isomorphic.

\begin{lemma} 
\label{lemma:isomorphic}
Let $\beta,\beta'$ be as above and let
$s:G\rightarrow \mu_n$ satisfy $s(1)=1$.
There is a canonical isomorphism
\[
\phi:\mathcal{H}_\beta(\overline{G}//B_k)_{gen}
\cong
\mathcal{H}_{\beta'}(\overline{G}//B_k)_{gen},
\]
as algebras. It sends $f(g,\zeta)=\zeta f(g,1)$ to
\[
\phi(f)(g,\zeta)=\zeta s(g) f(g,1).
\]
In other words, under this mapping, we have
\[
\phi(f_1*_{\beta'} f_2)(g,\zeta)
=(\phi(f_1)*_{\beta}\phi(f_2))(g,\zeta),
\]
where $*_\beta$ denotes the convolution
with respect to the $\beta$ cocycle
\[
(\phi_1*_\beta \phi_2)(g,\zeta)
=\zeta \int_{G}\phi_1(x,1)\phi_2((x^{-1}g,1))
\beta(x,x^{-1}g)^{-1} dx
\]
and, similarly, $*_{\beta'}$ denotes the convolution
with respect to the $\beta'$ cocycle.
\end{lemma}

{\bf proof}: The verification of this
is straightforward.      $\Box$

\section{The structure of $\mathcal{H}(G//B_k)$}
\label{section:sl2}

We want to determine the structure of the generalized Hecke algebra as
a finitely generated algebra with generators and relations.
Then we will do the same for the metaplectic analog and compare the
two.

The affine Weyl group $W_a$ of $G$ is generated by $s_1,s_2$ and
the Weyl group $W$ is generated by $s_1$. We will use the
Iwahori decomposition
\[
G=BW_aB
\]
to determine the structure of the generalized Hecke algebra.

Recall $(Ad g)(x)=gxg^{-1}$. Let $S=\{s_1,s_2,t_0,
t_0^{-1}\}\cup B$.

\begin{proposition} 
\label{prop:main}
The algebra $\mathcal{H}(G//B_k)$
is generated by the functions $f_g$, $g\in S$.
Assume that $vol(B_k)=1$.
These elements are subject to the relations below.

\noindent
(A) Relations for the $s_i$'s:

(i) $f_{s_1}*f_{s_1}=c\sum_x f_x$, for
$x\in ((-1)\cdot (Ad\, s_1)(B_k)\cdot B_k)/B_k$, where
$c=vol(B_k s_1 B_k)$, 

(ii) $f_{t_0}*f_{s_1}=f_{s_1}*f_{t_0^{-1}}$.

\noindent
(B) Relations for elements of $B$:

(i) $f_1$ is the identity of $\mathcal{H}(G//B_k)$,

(ii) $f_x*f_y=f_{xy}$, for $x,y\in B$,

(iii) $f_{xg}=f_{gx}=f_g$, for all $x\in B_k,\ g\in G$.

\noindent
(C) Mixed relations:

(i) $f_{s_i}*f_x=f_{Ad\, s_i(x)}*f_{s_i}$, for $x\in B\cap (Ad\, s_i)(B)$,

(ii) For $x\in \mathcal{O}^\times$,
\[
f_{s_1}*f_{u_+(x)}*f_{s_1}
=cf_{Ad\, s_1(u_+(x))}=cf_{u_-(-x)},
\]
where $c=vol(B_k s_i B_k)^2$,

(iii) For $x\in \mathcal{O}^\times$,
\[
f_{s_2}*f_{u_+(x)}*f_{s_2}
=cf_{Ad\, s_2(u_+(x))}=cf_{u_-(\pi x)},
\]
where $c=vol(B_k s_i B_k)^2$,

(iv) For all $g\in B$,
\[
\begin{aligned}
f_{t_0}*f_g*f_{t_0^{-1}}&=cf_{t_0gt_0^{-1}},\\
c&=\frac{vol(B_k t_0 B_k)^2}{vol(B_k t_0gt_0^{-1}B_k)}.
\end{aligned}
\]

\noindent
(The relations C(ii), C(iii) must be expressed in terms of the generators,
as in \cite{H}. The relation corresponding to C(ii) is
\[
\begin{aligned}
f_{s_1}*f_{u_+(x)}*f_{s_1}
&=c\bigl(f_{u_+(-x^{-1})}*f_{s_1}*f_{h(-x)}\\
&\qquad *f_{u_+(-x^{-1})}\bigr).
\end{aligned}
\]
The relation corresponding to C(iii) is similar.)

Furthermore, these relations form a defining set of relations for
$\mathcal{H}(G//B_k)$ as an algebra.
\end{proposition}

\begin{remark} 

{\it 
\begin{enumerate}

\item
This proposition has an analog for $G$ replaced by 
$GL(n)$ \cite{H}, Theorem 2.1, chapter 3. 

\item
Note that though $S$ is infinite, the set of
generators $\{f_g\ |\ g\in S\}$ is finite (and depends on
$k$), so the algebra is indeed finitely generated.

\item
 Possibly B(iii) and C(iii) are superfluous relations. 

\item
 Note that $s_1$ has order $4$, not order $2$ as the analogous
generator in \cite{H} does.

\item
 The constants are computed as follows:
\[
\begin{array}{c}
vol(B_k s_1 B_k)=
vol(B_k s_1 B_ks_1^{-1})\\
=vol(B_k s_1 B_ks_1^{-1})/vol(B_k)
=vol(U_k^- B_k/B_k).
\end{array}
\]
For $g\in B$, we have
\[
\begin{array}{c}
vol(B_k t_0gt_0^{-1}B_k)
=vol(B_k t_0B_k\cdot B_kgB_k\cdot B_kt_0^{-1}B_k)\\
=vol(B_k t_0B_k)vol(B_kgB_k)vol(B_kt_0^{-1}B_k)
=vol(B_k t_0B_k)^2.
\end{array}
\]
\end{enumerate}
} 
\end{remark}

Here are some general facts we will use in the proof below.

\begin{lemma} (a) $B_k$ is a normal subgroup of $B$.

(b) Let $U^+_k=\{u_+(x)\ |\ x\in \pi^k\mathcal{O}_F\}$
and $U^-_k=\{u_-(x)\ |\ x\in \pi^k\mathcal{O}_F\}$.
We have
\[
B_k\cdot (s_1^{-1}B_ks_1)
=(s_1^{-1}B_ks_1)\cdot B_k=U^-_kB_k,
\]
and $U^-_kB_k$ is a group with $B_k$ as a normal subgroup.

(c) In general, we have
\[
f_{g_1}*f_{g_2}(x)=\int_G f_{g_1}(y)f_{g_2}(y^{-1}x)dy
=vol(B_kg_1B_k\cap xB_kg_2^{-1}B_k).
\]

(d) In general, we have
\[
\begin{aligned}
\int_G f*g(x)\,dx
&=\int_G f(y)\,dy\int_G g(x)\,dx,\\
\int_G (f*g*h)(x)\,dx
&=\int_G f(x)\,dx\int_G g(y)\,dy\int_G h(z)\,dz.
\end{aligned}
\]
\label{lemma:9}
\end{lemma}

{\bf proof of lemma}: Parts (a) and (b) 
may be proven by direct matrix calculations. 
Parts (c), (d) are simple consequences of the definitions.
$\Box$

{\bf proof of proposition}: We shall follow the ideas in 
\cite{H}, proof of Theorem 2.1.

We begin by verifying A(i): $f_{s_1}*f_{s_1}=q\sum_{x\in X} f_x$, 
where $x$ runs over a complete set $X$ of representatives
of $((-1)\cdot (Ad\, s_1)(B_k)\cdot B_k)/B_k$. 

As in \cite{H}, parts (a), (b) of Lemma 
\ref{lemma:9} imply that 
\[
f_{s_1}*f_{s_1}=c\cdot \chi_{(-1)\cdot U^-_kB_k},
\]
where $\chi_S$ denotes the characteristic function of
a subset $S$ of $G$ and where $c\not= 0$ is some constant. 
In other words,
\[
f_{s_1}*f_{s_1}(g)=c\sum_{x\in X} f_x(g) ,\ \ \ g\in G.
\]
Plugging $g=1$ into both sides, we obtain
\[
vol(B_k s_1 B_k)=c.
\]

For part A(ii), note that part (c) of Lemma \ref{lemma:9}
immediately reduces the proof of A(ii) to the 
comparison of two volume integrals. 
The equality of these two expressions follow from the claim
$vol(B_kt_0^{-1}B_kg\cap B_ks_1B_k)
=vol(gB_kt_0B_k\cap B_ks_1B_k)$.
This is not trivial but the hard part is to show that 
$B_kt_0B_ks_1B_k=B_ks_1B_kt_0^{-1}B_k$. To see this, 
pick $b_1,b_2,b_3\in B_k$ and note 
\[
b_1t_0b_2s_1b_3=b_1t_0s_1u_1b_2'
=b_1s_1t_0^{-1}u_1b_2'
\]
for some $u_1\in U_k^-$ and $b_2'\in B_k$
(we've used part (a) of Lemma \ref{lemma:9}). A matrix calculation
shows that this is equal to $b_1s_1b_4t_0^{-1}b_2'$, for some
$b_4\in B_k$. This shows that $B_kt_0B_ks_1B_k\subset 
B_ks_1B_kt_0^{-1}B_k$. The other inclusion is proven similarly.
From this A(ii) follows without too much difficulty.

Part (c) of Lemma \ref{lemma:9} implies part B(i) of the 
proposition.

We now verify part B(ii). We know $vol(B_k)=1$ implies $vol(bB_k)=1$,
for each $b\in B$. 
This and Lemma \ref{lemma:general} implies part B(ii).

Part B(iii) follows from the definition of $f_g$.

We shall now verify C(i). Since $B_k$ is normal in $B$, we have
\[
vol(B_k s_i b B_k)=vol(B_k s_i b B_k b^{-1} b)=
vol(B_k s_i B_kb).
\]
Thus $vol(B_k s_i b B_k)=vol(B_k s_i B_k b)=vol(B_k s_i B_k)$.
This and Lemma \ref{lemma:general} implies 
$f_{s_i}*f_x=f_{s_ix}$, for $x\in B$. Likewise, if
$s_ixs_i^{-1}\in B$ then $f_{s_ixs_i^{-1}}*f_{s_i}=f_{s_ix}$.
This implies C(i).

We next verify C(ii). The proof of this part is sketched
in \cite{H}. The argument in \cite{H} gives
$f_{s_1}*f_{u_+(x)}*f_{s_1}=cf_{u_-(-x)}$, for
$x\in \mathcal{O}^\times$, for some constant $c\not= 0$.
Here $c$ may be evaluated using
part (d) of Lemma \ref{lemma:9}, which implies
\[
c={ vol(B_k s_i B_k)^2vol(B_k u_+(x) B_k)\over 
vol(B_k u_-(x) B_k)}=vol(B_k s_i B_k)^2.
\]
(Here we've used the fact that $u_+(x)\in B$, so
\[
vol(B_k u_+(x) B_k)=vol(B_k u_+(x))=vol(B_k)=1,
\]
and the fact that, for $k>0$, $s_1^{-1}B_ks_1$ is also a
normal subgroup of $B$, so
\[
\begin{aligned}
vol(B_k u_-(x) B_k)
&=vol(s_1^{-1}B_ks_1 u_+(-x) s_1^{-1}B_ks_1)\\
&=vol(s_1^{-1}B_ks_1 u_+(-x))\\
&=vol(s_1^{-1}B_ks_1)=1.
\end{aligned}
\]
)

The proof of C(iii) is very similar to part (ii), so
is omitted.

We shall now verify C(iv). We already know that, as a consequence of 
Lemma \ref{lemma:general}, we have $f_{t_0}*f_g=f_{t_0g}$
for all $g\in B$. This and part (c) of the above lemma implies that
\[
f_{t_0}*f_g*f_{t_0^{-1}}(x)=vol(B_kt_0gB_k\cap xB_kt_0B_k).
\]
This volume depends on $g$ and $x$, or more precisely on the
left cosets $gB_k$ and $xB_k$. In other words, we may
replace $g$ by any $gb$, $b\in B_k$, without changing the value
of the volume. The above volume is non-zero if and only if
$xb_1t_0b_2=b_3t_0gb_4$, for some $b_1,b_2,b_3,b_4$ in $B_k$.
By replacing $g$ by a suitable element in $gB_k$, we may assume
$x\in B_kt_0gt_0^{-1}B_k$. The support of $f_{t_0}*f_g*f_{t_0^{-1}}$
is therefore contained in (and thus must be equal to)
the support of $f_{t_0gt_0^{-1}}$. Therefore, there is a constant $c$
such that $f_{t_0}*f_g*f_{t_0^{-1}}=c\cdot f_{t_0gt_0^{-1}}$.
Part C(iv) follows from part (d) of the above lemma.

It remains to show that this list of relations is a 
defining set of relations for the algebra. To this end, 
we slightly modify the proof of \cite{H}, chapter 3, Theorem 2.1,
as follows.

For $g\in S$, let $\tilde{f}_g$ denote an abstract element.
Let $\tilde{\mathcal{H}}$ denote the free algebra generated by these 
elements satisfying the relations (A)-(C) in the 
proposition. We want to show that the obvious map
\begin{equation}
\tilde{\mathcal{H}}\rightarrow \mathcal{H}(G//B_k)
\label{eq:obvious}
\end{equation}
is an isomorphism. 

If $w\in W_a=<s_1,t_0>$ has the expression
$w=s_1^a t_0^b$, where $a\in \{0,1\}$ and $b$ is an integer,
then define $\tilde{f}_w=f_{s_1}^a*f_{t_0}^b$.
Let $\tilde{\mathcal{J}}$ denote the 
vector space span generated by the elements 
$\tilde{f}_g=\tilde{f}_x*\tilde{f}_w*\tilde{f}_y$, where
$g=xwy$, for $x,y\in B$ and $w \in W_a$. 
As in \cite{H}, pages 38-39, the relations (B,C) 
imply that $\tilde{f}_g$ depends only on $B_kgB_k$,
not on the particular representation $g=xwy$.
We claim that 
$\tilde{\mathcal{J}}$ is invariant under convolution by 
an abstract generator $\tilde{f}_g$, $g\in S$.
The proof of this invariance property, in \cite{H} pages 38-39,
uses a subset of the relations (or ones very similar) given in
(A)-(C) above. Therefore, it works in this case almost
verbatim and we omit the details. 
From this invariance property it follows that $\tilde{\mathcal{J}}$
contains the abstract generators. However, this
forces $\tilde{\mathcal{J}}$ to contain all of 
$\tilde{\mathcal{H}}$. We
conclude that each element of $\tilde{\mathcal{H}}$
may be written as a sum of elements of the form 
$\tilde{f}_g=\tilde{f}_x*\tilde{f}_w*\tilde{f}_y$, where
$g=xwy$. (One may think of this as a
``canonical form'' for the generators.)
This, using a relatively simple argument involving the
definitions (see \cite{H}), implies that the map (\ref{eq:obvious}) is 
injective.  

Since $B_k$ is normal in $B$, we have
\[
\begin{aligned}
vol(B_k b_1g b_2 B_k)
&=vol(b_1 b_1^{-1}B_k b_1 g b_2 B_k b_2^{-1} b_2)\\
&=vol(b_1 B_k g B_kb_2)=vol(B_k g B_k).
\end{aligned}
\]
This and Lemma \ref{lemma:general} implies 
$f_{b_1}*f_g*f_{b_2}=f_{b_1gb_2}$, for $g\in W^a$. 
Combining this with the Bruhat decomposition, we see
that $\mathcal{H}(G//B_k)$ is generated by elements of the form
$f_{b_1}*f_g*f_{b_2}$, for $g\in W^a$. 
Therefore, the map (\ref{eq:obvious}) is surjective.      $\Box$ 


\section{The structure of $\mathcal{H}(\overline{G}//B_k)_{gen}$}
\label{section:metaplectic}

This section is the metaplectic analog of the previous section.
We want to determine the structure of the metaplectic
generalized Hecke algebra as
a finitely generated algebra with generators and relations. 

We shall
always assume that $k$ is chosen so large 
that $\overline{B_k}$ splits.
By Lemma \ref{lemma:isomorphic},
we can (and do) choose the cocycle ${\beta'}$ defining
$\mathcal{H}(\overline{G}//B_k)_{gen}=
\mathcal{H}_{\beta'}(\overline{G}//B_k)_{gen}$
in such a way that the splitting of 
$\overline{B}_k$ is trivial.
If the residual characteristic of $F$ is relatively
prime to $2n$ then this condition is
vacuous since in that case, in fact,
$\overline{B}$ splits.

Let the Haar measure on $\overline{G}$ be normalized so that
$\overline{B_k}$ has measure one.
Convolution gives $\mathcal{H}(\overline{G}//B_k)_{gen}$ the structure of an
algebra. 

\begin{proposition} 
\label{prop:main2}
Assume that $vol(\overline{B}_k)=1$,
that $k$ is chosen so large that $\overline{B_k}$ splits,
and choose the splitting cocycle as above.
The algebra $\mathcal{H}(\overline{G}//B_k)_{gen}$
is generated by the functions $\overline{f}_g$, $g\in S$.
These elements are subject to the same relations 
as in Proposition \ref{prop:main}, except that $f_g$ must be
replaced by $\overline{f}_g$.
Furthermore, these relations form a defining set of relations for
$\mathcal{H}(\overline{G}//B_k)_{gen}$ as an algebra.
\end{proposition}

{\bf proof}: First we must verify that the claimed relations do 
indeed hold.

A(i): By Lemma \ref{lemma:convolve}, we have
\[
\overline{f}_{s_1}*\overline{f}_{s_1}(\overline{x})
=\int_{B_ks_1B_k}
\overline{f}_{s_1}(y^{-1}x)\beta'(y^{-1},x)\beta'(y,y^{-1})dy.
\]
Because of this, for basically 
the same reason as in the non-metaplectic case, 
$\overline{f}_{s_1}*\overline{f}_{s_1}$
is supported on $\overline{((-1)\cdot (Ad\, s_1)(B_k)\cdot B_k)}$.
The same reasoning as in the non-metaplectic case
implies A(i).

A(ii): Both $\overline{f}_{t_0}*\overline{f}_{s_1}$ and 
$\overline{f}_{s_1}*\overline{f}_{t_0^{-1}}$ are genuine functions which,
by the proof of part A(ii) in Proposition \ref{prop:main},
have the same support, in a neighborhood
of $(t_0s_1,\pm 1)$. We have, as in the proof of
A(ii) in the non-metaplectic case,
\[
\overline{f}_{t_0}*\overline{f}_{s_1}(g,1)
=\int_{\overline{B_ks_1B_k}}
\overline{f}_{t_0}((g,1)y^{-1})dy
=\int_{B_kt_0^{-1}B_kg\cap B_ks_1B_k} 
\beta'(y^{-1},g)dy
\]
and
\[
\overline{f}_{s_1}*\overline{f}_{t_0^{-1}}(g,1)
=\int_{gB_kt_0B_kg\cap B_ks_1B_k} \beta'(y^{-1},g)dy.
\]
We know that the ranges of these two integrals are 
equal,
by the proof in the non-metaplectic case above.
The validity of A(ii) follows.

B(i): For example, for $g\in \overline{G_s}$ we have
\[
\phi*\overline{f}_1(g)=\int_{\overline{G}}\phi(gx^{-1})\overline{f}_1(x)dx
=\int_{\overline{B_k}}\phi(gx^{-1})dx
=\phi(g).
\]
Similarly, $\overline{f}_1*\phi=\phi$, so $\overline{f}_1$ is an 
identity in $\mathcal{H}(\overline{G}//B_k)_{gen}$.

B(ii): We have
\[
\overline{f}_x*\overline{f}_y(g,1)
=\int_{B_kxB_k}
\overline{f}_{y}(h^{-1}g)
\beta'(h^{-1},g)\beta'(h,h^{-1})dh.
\]
This is non-zero if and only if $gy^{-1}=xb$ for some 
$b\in B_k$,
since $B_k$ is normal in $B$. But this implies
$supp(\overline{f}_x*\overline{f}_y)\subset 
\overline{B_kxyB_k}$. The claimed equality now follows.

B(iii): This is an immediate consequence of Lemma
\ref{lemma:genuine}.

C(i): The proof of this identity is analogous to the
proof in the non-metaplectic case (using Lemma
\ref{lemma:genuine} in place of Lemma \ref{lemma:general}
where appropriate) and is omitted.

C(ii), C(iii), C(iv): These are proven exactly as in
the non-metaplectic case so we only
sketch the argument. The idea is that
the proof in the non-metaplectic case shows that each of these
triple convolutions is supported on a single 
double $\overline{B}_k$-coset. On the other hand, the
support of each function in the Hecke algebra 
contains at least one double $\overline{B}_k$-coset.
since different double $\overline{B}_k$-cosets are
disjoint. Since convolutions preserve genuineness,
those functions whose support is precisely
one double $\overline{B}_k$-coset are the constant multiples of
the generators. This identity holds on all of
$\overline{G}$. The constant multiple can be computed as 
in the non-metaplectic case and since the measures
differ only by the measure on $\mu_n$,
the volume computations are similar. 

It remains to show that this list of relations is a 
defining set of relations for the algebra. This, too, is
exactly the same as the non-metaplectic case,
so is omitted.              $\Box$


\section{Application to representation theory}

The following is our main result:

\begin{theorem} 
\label{thrm:main}
If $k>0$ is such that $\overline{B_k}$ splits
then the map $f_{g}\longmapsto \overline{f}_{g}$, $g\in S$,
defines a bijection of algebras
$\mathcal{H}(G//B_k)\rightarrow \mathcal{H}(\overline{G}//B_k)_{gen}$.
\end{theorem}

{\bf proof}: This is an immediate consequence of
Propositions \ref{prop:main} and \ref{prop:main2}. $\Box$

\begin{corollary}
\label{corollary:main}
If $k>0$ is such that $\overline{B_k}$ splits then the bijection
$\mathcal{H}(G//B_k)\rightarrow \mathcal{H}(\overline{G}//B_k)_{gen}$
induces a set-theoretic 
bijection $\mathcal{H}(\overline{G}//B_k)_{gen}^\wedge
\rightarrow \mathcal{H}(G//B_k)^\wedge$ on the sets of equivalence
classes of irreducible (genuine, in the case of $\overline{G}$)
representations. Furthermore, if
\[
\begin{array}{ccc}
\mathcal{H}(G//B_k)&\rightarrow &\mathcal{H}(\overline{G}//B_k)_{gen}\\
f & \longmapsto &\overline{f}
\end{array}
\]
and 
\[
\begin{array}{ccc}
\mathcal{H}(\overline{G}//B_k)_{gen}^\wedge 
  &\rightarrow &\mathcal{H}(G//B_k)^\wedge\\
\overline{\pi}&\longmapsto &\pi
\end{array}
\]
then ${\rm tr}(\overline{\pi}(\overline{f}))={\rm tr}(\pi(f))$.
\end{corollary}

{\bf proof}: Let
$\eta:\mathcal{H}(G//B_k)\rightarrow \mathcal{H}(\overline{G}//B_k)_{gen}$
and let $\eta^*:\mathcal{H}(\overline{G}//B_k)_{gen}^\wedge
\rightarrow \mathcal{H}(G//B_k)^\wedge$. Suppose
$\overline{\pi}(\overline{f})=\overline{\pi}'(\overline{f})$,
for all $\overline{f}\in \mathcal{H}(\overline{G}//B_k)_{gen}$. If 
$\overline{f}=\phi(f)$, for some $f\in \mathcal{H}(G//B_k)$,
and if $\pi=\eta^*(\overline{\pi})$ and $\pi'=\eta^*(\overline{\pi}')$, 
for some $\pi,\pi'\in \mathcal{H}(G//B_k)^\wedge$, then 
$\pi(f)=\pi'(f)$, for all $f\in \mathcal{H}(G//B_k)$.
This implies $\pi=\pi'$, so $\eta^*$ is injective.

The equality of the traces is by definition of $\eta$. $\Box$ 

\begin{definition} 
Let $K$ denote a compact open subgroup of $\overline{G}$.
A locally constant
complex-valued function $f$ on $\overline{G}$
is called {\bf left} (resp., {\bf right}) {\bf Hecke finite}
if the vector space spanned by all functions of the form
$h*f$ (resp., $f*h$), for $h\in \mathcal{H}(\overline{G}//K)$,
is finite dimensional. Let $\mathcal{A}_{gen}(\overline{G})$
denote the space of all locally constant
complex-valued genuine functions
$f$ on $\overline{G}$ for which
\begin{itemize}
\item
$f$ is double $K$-finite,
\item
$f$ is right- and left-Hecke finite. 
\end{itemize}
This space is called the space of {\bf matrix
coefficients} of $\overline{G}$.
If $(\pi,V)$ is a genuine admissible representation
of $\overline{G}$ then the span of the
elements $f(g)=<v_0^*,\pi(g)v_0>$ ($g\in \overline{G}$),
for all $v_0\in V$ and $v_0^*\in V^*$,
the contragredient of $V$, is called the space 
of {\bf matrix coefficients of $\pi$} and is
denoted $\mathcal{A}_{gen}(\pi)$.
\end{definition}

Let $\rho$ (resp, $\lambda$) denote the right 
(resp., left) regular representation of $\overline{G}$ 
on $\mathcal{A}_{gen}(\overline{G})$.
The next proposition establishes that the metaplectic 
analogs of (the relevant parts of) Corollary 1.10.6 in
\cite{Sil} for $\overline{G}$ are true. The result below
implies that we can ``explicitly'' realize the space of any genuine
irreducible representation of $\overline{G}$ as a subspace
of $\mathcal{A}_{gen}(\overline{G})$, a fact we will make use of
later.

\begin{proposition} 
\label{prop:silberger}
Let $f\in \mathcal{A}_{gen}(\overline{G})$
and let $V_f$ denote the $\overline{G}$-module
spanned by the translates $\rho(g)f$,
$g\in \overline{G}$, and let $\pi$ denote
the restriction of $\rho$ to $V_f$. We have
\begin{itemize}
\item[(a)]
There exist elements $v_0\in V_f$ and $v_0^*\in V_f^*$,
the contragredient of $V_f$, such that
$f(g)=<v_0^*,\pi(g)v_0>$, for
all $g\in \overline{G}$.
\item[(b)]
$\mathcal{A}_{gen}(\pi)$ is spanned by 
$\lambda(x)\rho(y)f$, $x,y\in \overline{G}$.
\end{itemize}
\end{proposition}

{\bf proof}:
The proof of this follows from the general
arguments given in \cite{Sil}. The results in that
part of the book \cite{Sil} hold for more
general totally disconnected groups,
so one need only check that the
argument remains valid when ``genuineness'' 
is also assumed.
$\Box$

We need to recall some facts regarding the relationship between 
smooth representations of $\overline{G}$ and finite dimensional
$\mathcal{H}(\overline{G}//B_k)_{gen}$-modules. To this end,
we recall briefly how one can construct 
smooth representations of $\overline{G}$ from 
finite dimensional $\mathcal{H}(\overline{G}//B_k)_{gen}$-modules -
both the so-called ``induced'' and ``produced'' modules 
in \S 2 of Borel \cite{Bo} will do this job. 
The ``inverse'' construction is then given along with
some of the basic properties of these constructions
in Proposition \ref{prop:borel} below.

\begin{definition}
If $(r,W)$ is a finite dimensional 
$\mathcal{H}(\overline{G}//B_k)$-module then we define 
$(I(r),I(W))$ to be the {\bf induced} $\overline{G}$-module
and $(P(r),P(W))$ to be the (smooth) {\bf produced} $\overline{G}$-module
constructed in \cite{Bo}. Analogously, if $(r,W)$ is a finite dimensional 
$\mathcal{H}(\overline{G}//B_k)_{gen}$-module then we define 
$(I_{gen}(r),I_{gen}(W))$ to be the {\bf induced} genuine 
$\overline{G}$-module and 
$(P_{gen}(r),P_{gen}(W))$ to be the {\bf produced} genuine 
$\overline{G}$-module obtained by replacing 
$C_c(\overline{G}//B_k)$ by $C_c(\overline{G}//B_k)_{gen}$
in the contructions in \cite{Bo}.
\end{definition}

\begin{proposition}
\label{prop:borel}
Let $(r,W)$ be a finite dimensional
$\mathcal{H}(\overline{G}//B_k)_{gen}$-module and let
$(\pi,V)$ be a genuine smooth representation of $\overline{G}$.
\begin{itemize}
\item[(a)] 
There are natural $\mathcal{H}(\overline{G}//B_k)_{gen}$-module
isomorphisms
\[
W\cong I_{gen}(W)^{B_k}\ \ 
{\rm and}\ \ W\cong P_{gen}(W)^{B_k}. 
\]
\item[(b)] 
$I_{gen}(W)$ is generated, as a $\overline{G}$-module,
by $I_{gen}(W)^{B_k}$.
\item[(c)]
The natural restriction maps yield isomorphisms
\[
\rho_P : {\rm Hom}_{\overline{G}}(V,P_{gen}(W))
\rightarrow 
 {\rm Hom}_{\mathcal{H}(\overline{G}//B_k)_{gen}}(V^{B_k},W).
\]
and
\[
\rho_I : {\rm Hom}_{\overline{G}}(I_{gen}(W),V)
\rightarrow 
 {\rm Hom}_{\mathcal{H}(\overline{G}//B_k)_{gen}}(W,V^{B_k}).
\]
\item[(d)]
The $\mathcal{H}(\overline{G}//B_k)_{gen}$-module $(r,W)$
is irreducible if and only if its induced $\overline{G}$-module
$(I_{gen}(r),I_{gen}(W))$ is irreducible. Moreover, for an
irreducible $\mathcal{H}(\overline{G}//B_k)_{gen}$-module $(r,W)$,
we have $I_{gen}(r)\cong P_{gen}(r)$.
\end{itemize}
\end{proposition}

{\bf proof}: In the case $n=1$,
the statements in the proposition above are either in \S\S 1-2 \cite{Bo} 
as stated or follow from Borel's results using Schur's lemma
(see for example, Lemma 3.14 in \cite{CR}).
In general, one need only check that the
argument remains valid when ``genuineness'' is also assumed.
$\Box$

\begin{definition} 
We define a genuine
admissible irreducible representation
of $\overline{G}$ to be {\bf supercuspidal} 
if and only if its matrix coefficients are compactly supported.
We define a genuine
admissible irreducible representation
of $\overline{G}$ to be {\bf square-integrable} 
if and only if
its matrix coefficients belong to $L^2(\overline{G})$.

We define an irreducible finite dimensional
$\mathcal{H}(\overline{G}//B_k)_{gen}$-module $(r,W)$ to be 
{\bf supercuspidal} (resp., {\bf square-integrable})
if and only if $I_{gen}(r)$ is a supercuspidal 
(resp., square-integrable) representation of
$\overline{G}$. 
\end{definition} 

Next we observe that, as in the case
of $GL(n)$, the correspondence of Corollary \ref{corollary:main}
preserves supercuspidals and discrete series 
representations (\cite{H}, page 33).

\begin{corollary}
Assume $k>0$ is as in Corollary \ref{corollary:main}.
\begin{itemize}
\item An admissible genuine representation
$\overline{\pi}\in
\mathcal{H}(\overline{G}//B_k)_{gen}^\wedge$ 
is supercuspidal if and only if
$\pi$ is, where $\pi,\overline{\pi}$ are as in Corollary
\ref{corollary:main}.

\item
$\overline{\pi}$ is square-integrable 
if and only if $\pi$ is, where $\pi,\overline{\pi}$ 
are as in Corollary \ref{corollary:main}.
\end{itemize}
\end{corollary}

{\bf proof}: Suppose that the mapping
$\mathcal{H}(\overline{G}//B_k)_{gen}^\wedge
\rightarrow \mathcal{H}(G//B_k)^\wedge$
sends a representation $(\overline{r},\overline{W})$ 
to a representation $(r,W)$. Note that the construction of
correspondence in Corollary \ref{corollary:main} implies that
$W=\overline{W}$.
From the construction of $I_{gen}(\overline{W})=
C_c(\overline{G}/B_k)_{gen}
\otimes_{\mathcal{H}(\overline{G}//B_k)_{gen}}\overline{W}$
(more precisely from the definition of 
$\otimes_{\mathcal{H}(\overline{G}//B_k)_{gen}}$), we
see that the $r(\overline{f}_g)$-action on 
$\overline{W}\cong I_{gen}(\overline{W})^{B_k}$
corresponds to the $\overline{f}_g$-action (via convolution)
on $C_c(\overline{G}/B_k)_{gen}$. 
Identifying $\overline{W}$ with $I_{gen}(\overline{W})^{B_k}$, 
we see that each element of
$\overline{W}$ is Hecke-finite and may therefore be regarded as 
a matrix coefficient of
$\overline{G}$. Let $(\overline{\pi},\overline{V})$
denote the $\overline{G}$-module generated by
an element $w\in \overline{W}\cong I_{gen}(\overline{W})^{B_k}$. 
By Proposition \ref{prop:silberger} (or Proposition \ref{prop:borel} (b)), 
$\overline{\pi}=I_{gen}(\overline{r})$.
%Define the support of an element
%$w\in \overline{W}$ to be the set of $g\in \overline{G}$ such that 
%$r(\overline{f}_g)(w)\not= 0$. 
From this and the fact that $W=\overline{W}$, it 
follows that if the elements of $\overline{W}$
are compactly supported (when regarded as matrix coefficients of $\overline{G}$)
then the elements of $W$ are also are compactly supported
(when regarded as matrix coefficients of $G$).
Now assume that $(\overline{r},\overline{W})$ is
supercuspidal, so the matrix coefficients of $I_{gen}(\overline{r})$
have compact support. As indicated above, $(r,W)$ is supercuspidal as well. 
This established the first part of the corollary.

The second part of the corollary is proved similarly,
hence is omitted.
$\Box$

\section{The Iwahori algebra}
\label{section:savin}

In the case where $k=0$ and ${\rm gcd}(p,2n)=1$, 
our result is a special
case of a general isomorphism between Iwahori algebras,
due to Savin \cite{S}, which we recall:

\begin{theorem} (Savin) Assume ${\rm gcd}(p,2n)=1$.
There is an isomorphism
of algebras
$\mathcal{H}(\overline{G}//B)_{gen}\rightarrow \mathcal{H}(G//B)$.
\end{theorem}

\begin{corollary}
 Assume ${\rm gcd}(p,2n)=1$. The isomorphism
$\mathcal{H}(\overline{G}//B)_{gen}\rightarrow \mathcal{H}(G//B)$
induces a set-theoretic 
bijection $\mathcal{H}(\overline{G}//B)_{gen}^\wedge
\rightarrow \mathcal{H}(G//B)^\wedge$ on the sets of equivalence
classes of irreducible (genuine, in the case of $\overline{G}$)
representations with a $B$-fixed vector.
Furthermore, if
\[
\begin{array}{ccc}
\mathcal{H}(G//B)&\rightarrow &\mathcal{H}(\overline{G}//B)_{gen}\\
f & \longmapsto &\overline{f}
\end{array}
\]
and 
\[
\begin{array}{ccc}
\mathcal{H}(\overline{G}//B)_{gen}^\wedge &\rightarrow &\mathcal{H}(G//B)^\wedge\\
\overline{\pi}&\longmapsto &\pi
\end{array}
\]
then ${\rm tr}(\overline{\pi}(\overline{f}))={\rm tr}(\pi(f))$.
\end{corollary}

%\QTP{bibsection}

\begin{thebibliography}{99}

\bibitem[Bo]{Bo} A. Borel, ``Admissible representations of 
a semi-simple group
over a local field with vectors fixed under an Iwahori 
subgroup,'' Invent. Math. \underline{35}(1976)233-259

\bibitem[CR]{CR} C. Curtis, I. Reiner, {\it Methods 
of representation theory, volume 1}, Wiley Interscience, 
John Wiley and Sons,
New York, 1981

\bibitem[F]{F} Y. Flicker, ``Automorphic forms on covering groups of 
$GL(2)$'', Inv. Math. \underline{57}(1980)119-182

\bibitem[FK]{FK} ------ and D. Kazhdan, ``Metaplectic 
correspondence'', Publ. Math. IHES \underline{64}(1986)53-110

\bibitem[G]{G} S. Gelbart, {\it Weil's representation and 
the spectrum of the metaplectic group},
SLN Math \underline{530}, 1976

\bibitem[H]{H} R. Howe (with the collaboration of A. Moy),
{\it Harish-Chandra homomorphisms for p-adic groups}, 
CBMS \underbar{59}, AMS (1985).

%\bibitem[J]{J} D. Joyner, \underbar{A}
%\underbar{application}
%\underbar{of} \underbar{Arthur's}
%\underbar{invariant} \underbar{trace}
%\underbar{formula}
%\underbar{for} \underbar{$SL(2)$}, preprint

\bibitem[KP]{KP} D. Kazhdan and  S. Patterson, ``Metaplectic forms'', 
Publ. Math. IHES \underline{59}(1984)35-142

\bibitem[K]{K} T. Kubota, ``Topological covering of $SL(2)$ over a 
local field'', J.
Math. Soc. Japan \underline{19} (1967)114-121

%\bibitem[Ma]{Ma} D. Manderscheid, ``On the supercuspidal 
%representations of
%$SL(2)$ and its 2-fold cover, I'', Math. Ann. 
%\underbar{266}(1987)287-295

\bibitem[M]{M} L. Morris, ``P-cuspidal representations'', 
Proc. London Math. Soc. \underline{57}(1988)329-356

\bibitem[RS]{RS} S. Rallis and G. Schiffmann,
``Repr\'esentations supercuspidales du
groupe m\'etaplectique,'' J. Math. Kyoto
Univ. \underbar{17}(1977)567-603

\bibitem[S]{S} G. Savin, ``A local Shimura correspondence'',
Math Ann. \underbar{280}(1988)185-190

\bibitem[Sch]{Sch} J. Schultz, ``Lifting of characters of 
$\overline{SL_2(F)}$ and
$SO_{2,1}(F)$, for $F$ nonarchimedean'',
PhD thesis (supervised by J. Adams),
Univ. Maryland at College Park, expected summer 1998

\bibitem[Sil]{Sil} A. Silberger, {\it Introduction to
harmonic analysis on reductive p-adic groups}, Princeton Univ. 
Press, Princeton, 1979

\bibitem[W]{W} J.-L. Waldspurger, ``Correspondence de Shimura'',
J. Math. Pure Appl. \underline{59}(1980)1-133
\end{thebibliography}

\end{document}
